% Mục 3.2 (bản rút gọn). Rút từ sections/impl_barrier.tex. Nguồn: include/barrier.hpp.
Barrier cho $n$ thread bảo đảm không thread nào ra khỏi lần gọi \texttt{wait()} thứ $k$ trước khi cả $n$
thread đã gọi lần thứ $k$, và dùng lại được qua mọi episode. Barrier một cờ ``mở/đóng'' không dùng lại
được. Sense-reversing barrier~\cite{hensgen1988barrier} không đặt lại cờ mà đảo nó mỗi episode. Mỗi
thread chờ cờ bằng \emph{local sense} của mình, nên hai episode kề nhau chờ hai giá trị khác nhau. Mỗi
thread tới barrier tăng bộ đếm \texttt{count\_} bằng \texttt{fetch\_add}. Thread tới cuối đặt lại bộ đếm
về 0 rồi ghi local sense của nó vào cờ chung \texttt{sense\_}. Các thread khác chờ tới khi cờ bằng local
sense của mình. Sau đó mỗi thread đảo local sense của mình cho episode sau.
Thread tới cuối phải đặt lại bộ đếm trước khi đảo cờ. Nếu đảo trước, thread vừa được thả có thể tới
episode sau, rồi bị phép đặt lại muộn xoá mất lần tới, và mọi thread chờ mãi.

\texttt{fetch\_add} cần acq\_rel (\texttt{ldaddal}). Lý do nằm ở hai khái niệm của mô hình bộ nhớ C++.
Một phép ghi release và một phép đọc acquire đọc được giá trị của nó tạo quan hệ \emph{happens-before}:
mọi thứ một thread ghi trước phép ghi release, thread kia nhìn thấy sau phép đọc acquire. Một phép ghi
release còn mở đầu một \emph{release sequence}, là dãy gồm nó và mọi phép đọc-sửa-ghi atomic đứng sau
nó trên cùng biến. Phép đọc acquire đọc được bất kỳ phần tử nào của dãy đều đồng bộ với phép ghi đầu
dãy. \texttt{fetch\_add} là phép đọc-sửa-ghi, nên các \texttt{fetch\_add} lên \texttt{count\_} nối nhau
thành một dãy như vậy. Vì thế \texttt{fetch\_add} của thread tới cuối đồng bộ với \texttt{fetch\_add}
của mọi thread tới trước, và thread tới cuối thấy mọi phép ghi của chúng trước barrier. Phép ghi release
lên \texttt{sense\_} rồi phép đọc acquire của thread chờ chuyển tiếp các phép ghi đó cho mọi thread đang
chờ.

\texttt{count\_}, \texttt{sense\_} và mỗi local sense nằm trên các cache line riêng. Local sense nằm trong barrier, đánh chỉ số bằng \texttt{tid}, vì một biến \texttt{thread\_local}
sẽ bị hai barrier dùng xen kẽ làm hỏng.

Các cách chờ gồm quay (\PCname{sr}), ba chính sách backoff như lock, và chờ chặn (\PCname{sr_block}) bằng
\texttt{std::atomic::wait} và \texttt{notify\_all}. Theo mã nguồn libc++ và libstdc++, \texttt{wait}
quay một lúc rồi mới chặn thật. Các barrier thư viện gồm \texttt{std::barrier}, barrier dựng từ mutex và
condition variable (\PCname{condvar}), \texttt{pthread\_barrier\_t} (Linux) và \texttt{omp barrier}.
